\chapter{Physical systems and models}
\label{chap:systems}
%=============================================================================

The three preceding chapters built a formalism.  This one puts it to work on a
handful of systems.

The models collected here share a property that makes them invaluable: each is
simple enough to be solved exactly, at least for small numbers of particles,
yet each contains a genuine many-body mechanism.  The Lipkin model has a
collective mode driven by a two-body interaction that scatters pairs across a
gap; the pairing model has superfluid correlations, and a small pair-breaking
addition to it turns the particle-hole channel on and gives us the model that
carries the whole of Chapter~\ref{chap:mfapplications}; the Fermi-Hubbard
model has the competition between kinetic delocalisation and on-site repulsion
that drives the Mott transition; its own strong-coupling limit is the quantum
Heisenberg model, the simplest interacting spin system there is; and the
Calogero-Sutherland model is exactly solvable for arbitrary particle number,
with a ground state that is a pure Jastrow function.  Because the exact answer
is available, these models are the benchmarks against which every approximate
method in the rest of this book is measured.

They also share a structure.  Every one of them is written as
\begin{equation}
  \hat{H} = \sum_{pq}\langle p|\hat{h}_0|q\rangle a_p^{\dagger}a_q
    + \frac{1}{4}\sum_{pqrs}\langle pq|\hat{v}|rs\rangle_{\mathrm{AS}}
      a_p^{\dagger}a_q^{\dagger}a_s a_r ,
  \label{eq:4-generalH}
\end{equation}
the general second-quantised Hamiltonian of
Chapter~\ref{chap:secondquant}, with a particular and very restricted choice
of matrix elements.  What differs between them is which single-particle
states there are and which matrix elements are allowed to be non-zero.  The
symmetries that follow from those restrictions are what make the models
tractable, and identifying them is the first step in every case.  The one
exception is the Heisenberg model of Section~\ref{sec:heisenberg}, which is
written in spin operators rather than in fermion operators; the Jordan-Wigner
transformation of Section~\ref{sec:pauli} maps it back onto
Eq.~(\ref{eq:4-generalH}), and the mapping is instructive in its own right.

Throughout, the detailed algebra is relegated to the exercises at the end of
the chapter, where full solutions are given.  The main text states each model,
identifies its symmetries, and reports what the exact solution looks like.
All numbers quoted are produced by \texttt{models.py}.

%----------------------------------------------------------------------
\section{The Lipkin model}
\label{sec:lipkin}

\index{Lipkin model}
The Lipkin-Meshkov-Glick model~\cite{lipkin1965} describes $N$ fermions
distributed over two levels, each of degeneracy $\Omega$.  The levels carry
the quantum number $\sigma=\pm1$: the upper has $\sigma=+1$ and energy
$+\varepsilon/2$, the lower has $\sigma=-1$ and energy $-\varepsilon/2$.  The
substates of each level are labelled $p=1,\dots,\Omega$, so a single-particle
state is $\bket{p\sigma}=a_{p\sigma}^{\dagger}\bket{0}$.

The Hamiltonian is
\begin{equation}
  \hat{H} = \hat{H}_0 + \hat{H}_1 + \hat{H}_2,
  \label{eq:4-lipkinH}
\end{equation}
with
\begin{align}
  \hat{H}_0 &= \frac{1}{2}\varepsilon\sum_{p\sigma}\sigma\,
    a_{p\sigma}^{\dagger}a_{p\sigma},
  \label{eq:4-lipkinH0}\\
  \hat{H}_1 &= \frac{1}{2}V\sum_{p p'\sigma}
    a_{p\sigma}^{\dagger}a_{p'\sigma}^{\dagger}
    a_{p'-\sigma}a_{p-\sigma},
  \label{eq:4-lipkinH1}\\
  \hat{H}_2 &= \frac{1}{2}W\sum_{p p'\sigma}
    a_{p\sigma}^{\dagger}a_{p'-\sigma}^{\dagger}
    a_{p'\sigma}a_{p-\sigma}.
  \label{eq:4-lipkinH2}
\end{align}
The operator $\hat{H}_1$ moves a \emph{pair} of particles from one level to
the other, with strength $V$; $\hat{H}_2$ is a spin-exchange term of strength
$W$, which moves a pair from $(p\sigma,p'-\sigma)$ to $(p-\sigma,p'\sigma)$.
Neither changes the substate labels $p$, and this is the crucial restriction:
the interaction is completely degenerate in $p$.

\paragraph{Quasispin.}
The degeneracy in $p$ means the model has an $SU(2)$ symmetry.  Define the
\emph{quasispin} operators
\begin{equation}
  \hat{J}_{\pm} = \sum_p a_{p\pm}^{\dagger}a_{p\mp},
  \qquad
  \hat{J}_z = \frac{1}{2}\sum_{p\sigma}\sigma
    a_{p\sigma}^{\dagger}a_{p\sigma},
  \qquad
  \hat{J}^2 = \hat{J}_+\hat{J}_- + \hat{J}_z^2 - \hat{J}_z .
  \label{eq:4-quasispin}
\end{equation}
Exercise~\ref{ex:4-lipkin} shows, using only the anticommutation relations of
Chapter~\ref{chap:secondquant}, that these satisfy the angular-momentum
algebra
\begin{equation}
  [\hat{J}_z,\hat{J}_{\pm}] = \pm\hat{J}_{\pm},
  \qquad
  [\hat{J}_+,\hat{J}_-] = 2\hat{J}_z,
  \qquad
  [\hat{J}^2,\hat{J}_{\pm}] = [\hat{J}^2,\hat{J}_z] = 0,
  \label{eq:4-su2}
\end{equation}
and that the Hamiltonian collapses to
\begin{equation}
  \boxed{\;
  \hat{H} = \varepsilon\hat{J}_z
    + \frac{1}{2}V\left(\hat{J}_+^2+\hat{J}_-^2\right)
    + \frac{1}{2}W\left(-\hat{N} + \hat{J}_+\hat{J}_- + \hat{J}_-\hat{J}_+
      \right) \;}
  \label{eq:4-lipkinquasispin}
\end{equation}
with $\hat{N}=\sum_{p\sigma}a_{p\sigma}^{\dagger}a_{p\sigma}$ the number
operator.  Every trace of the substate index $p$ has disappeared.

The consequence is decisive.  Since $[\hat{H},\hat{J}^2]=0$ and
$[\hat{H},\hat{N}]=0$, both $J$ and $N$ are good quantum numbers, and the
Hamiltonian is block diagonal in $J$.  A Hilbert space of dimension
$\binom{2\Omega}{N}$ collapses into blocks of dimension $2J+1$.  For $N=4$
particles at half filling the largest block is $J=2$, of dimension five: a
$70$-dimensional problem has become a $5\times5$ matrix.

\paragraph{The $J=2$ block.}
The states are labelled $\bket{J,J_z}$ with $J_z=-J,\dots,J$.  The
lowest, with all four spins down, is
\begin{equation}
  \bket{2,-2} = a_{1-}^{\dagger}a_{2-}^{\dagger}
                a_{3-}^{\dagger}a_{4-}^{\dagger}\bket{0},
  \label{eq:4-lipkinlowest}
\end{equation}
and the remaining four follow by applying $\hat{J}_+$ with the standard
normalisations
\begin{equation}
  \hat{J}_{\pm}\bket{J,J_z}
   = \sqrt{J(J+1)-J_z(J_z\pm1)}\,\bket{J,J_z\pm1} .
  \label{eq:4-ladder}
\end{equation}
Working out the matrix elements of
Eq.~(\ref{eq:4-lipkinquasispin}) -- again Exercise~\ref{ex:4-lipkin} -- gives
for $\varepsilon=2$, $V=-1/3$, $W=-1/4$
\begin{equation}
  \bm{H}_{J=2}^{(1)} =
  \begin{bmatrix}
    -4 & 0 & -\sqrt{6}/3 & 0 & 0\\
     0 & -11/4 & 0 & -1 & 0\\
    -\sqrt{6}/3 & 0 & -1 & 0 & -\sqrt{6}/3\\
     0 & -1 & 0 & 5/4 & 0\\
     0 & 0 & -\sqrt{6}/3 & 0 & 4
  \end{bmatrix},
  \label{eq:4-lipkinmatrix1}
\end{equation}
whose eigenvalues are $-4.21288$, $-2.98607$, $-0.91914$, $1.48607$ and
$4.13201$.  The ground state is
\begin{equation}
  \bket{\psi_0} = 0.96735\bket{2,-2} + 0.25221\bket{2,0}
                + 0.02507\bket{2,2},
  \qquad E_0 = -4.21288 .
  \label{eq:4-lipkinground1}
\end{equation}
Note the structure: the interaction connects only states differing by two
units of $J_z$, because $\hat{J}_{\pm}^2$ moves a pair.  The matrix therefore
splits further into even and odd $J_z$.

Increasing the interaction to $V=-4/3$, $W=-1$ gives eigenvalues $-7.75122$,
$-7.47214$, $-1.55581$, $1.47214$, $5.30704$ and the ground state
\begin{equation}
  \bket{\psi_0} = 0.64268\bket{2,-2} + 0.73816\bket{2,0}
                + 0.20515\bket{2,2},
  \qquad E_0 = -7.75122 .
  \label{eq:4-lipkinground2}
\end{equation}
In the first case the unperturbed configuration $\bket{2,-2}$ carries $94\%$
of the probability and a single Slater determinant is a good approximation.
In the second the interaction matrix elements are comparable to the
single-particle spacing, $\bket{2,0}$ has become the dominant component, and
no single determinant will do.  This is the same story the occupation numbers
of Section~\ref{sec:naturalorbitals} told, in a model small enough to see the
whole of it.

\begin{notebox}
\textbf{Why this model matters.} The Lipkin model was invented to test
many-body methods, and it still is.  Because the exact answer is a
$5\times5$ eigenvalue problem while the interaction can be tuned from weak to
strong, it exposes exactly where an approximation breaks down: Hartree-Fock
theory undergoes a spurious phase transition at a critical coupling, the
random-phase approximation predicts an excitation energy that goes imaginary
there, and coupled-cluster theory tracks the exact result much further.  Every
one of these statements can be checked against
Eq.~(\ref{eq:4-lipkinground1}) and its analogues.
\end{notebox}

\paragraph{Mapping to qubits.}
The Lipkin model is a natural first target for quantum algorithms, and how one
maps it onto qubits is worth a moment.  There are three obvious schemes.

The \emph{occupation mapping} assigns one qubit to each single-particle state
$(p,\sigma)$, the qubit being $\bket{1}$ or $\bket{0}$ according as that state
is occupied or empty.  It costs $2\Omega$ qubits.

The \emph{level mapping} exploits the fact that there are only two levels.
Restricting to half filling, $N=\Omega$, each doublet $\left((p,+1),
(p,-1)\right)$ holds exactly one particle, so one qubit per doublet suffices:
$\bket{0}$ for the particle in the upper state and $\bket{1}$ for the lower.
This costs $\Omega$ qubits, half of the occupation mapping.  It has a second
advantage: moving a pair changes the state of two doublets, so any operator in
the Hamiltonian is at most a two-qubit gate.

The \emph{spin mapping} goes further and assigns a qubit to each state
$\bket{J,J_z}$.  For a given $J$ there are only $2J+1$ of these, so with
$J_{\max}=N/2$ the largest register is $N+1$ qubits.  The price is that one
must run a separate circuit for each $J$ and take the lowest of the resulting
energies.

Under the level mapping the quasispin operators become sums of Pauli matrices.
Writing $\hat{J}_z=\sum_p j_z^{(p)}$ and $\hat{J}_{\pm}=\sum_p
j_{\pm}^{(p)}$ with
\begin{equation}
  j_z^{(p)} = \frac{1}{2}\sum_{\sigma}\sigma
    a_{p\sigma}^{\dagger}a_{p\sigma},
  \qquad
  j_{\pm}^{(p)} = a_{p\pm}^{\dagger}a_{p\mp},
  \label{eq:4-onebodyquasispin}
\end{equation}
each doublet operator is a single-qubit operator, and the correspondence with
the Pauli matrices of Section~\ref{sec:pauli} is
\begin{equation}
  j_z^{(p)} \;\longleftrightarrow\; \tfrac12 Z_p,
  \qquad
  j_{\pm}^{(p)} \;\longleftrightarrow\; \tfrac12\left(X_p \pm \imath Y_p\right).
  \label{eq:4-paulimapping}
\end{equation}
The Hamiltonian~(\ref{eq:4-lipkinquasispin}) therefore becomes a sum of Pauli
strings of weight at most two -- exactly the form a variational quantum
eigensolver requires.  The circuits themselves belong to the
quantum-computing part of this book; what matters here is that the mapping
costs $\Omega$ qubits and two-qubit interactions only, which is a direct
consequence of the $SU(2)$ structure identified above.

%----------------------------------------------------------------------
\section{The pairing model}
\label{sec:pairingmodel}

\index{pairing model}
The pairing model has appeared repeatedly already -- as the test case for the
Lanczos algorithm in Section~\ref{sec:lanczos}, for the Schmidt decomposition
in Section~\ref{sec:schmidt}, and as the three-level exercise of week~35.  We
now set it out properly.

The single-particle space consists of $L$ doubly degenerate, equally spaced
levels labelled $p=1,2,\dots,L$ with spin $\sigma=\pm1$.  The Hamiltonian is
\begin{equation}
  \hat{H} = \hat{H}_0 + \hat{V},
  \qquad
  \hat{H}_0 = \xi\sum_{p\sigma}(p-1)\,a_{p\sigma}^{\dagger}a_{p\sigma},
  \qquad
  \hat{V} = -\frac{1}{2}g\sum_{pq}
    a^{\dagger}_{p+}a^{\dagger}_{p-}a_{q-}a_{q+} .
  \label{eq:4-pairingH}
\end{equation}
The spacing $\xi$ is set to unity without loss of generality.  The interaction
has a single term and a single strength $g$: it destroys a pair of particles
in level $q$ with opposite spins and creates one in level $p$.  Nothing else
happens.  In the language of Eq.~(\ref{eq:4-generalH}), only the matrix
elements $\langle p+\,p-|\hat{v}|q+\,q-\rangle$ are non-zero, and they are all
equal.

\paragraph{Symmetries.}
Exercise~\ref{ex:4-pairing} shows that both $\hat{H}_0$ and $\hat{V}$ commute
with the spin projection and the total spin,
\begin{equation}
  \hat{S}_z = \frac{1}{2}\sum_{p\sigma}\sigma
    a^{\dagger}_{p\sigma}a_{p\sigma},
  \qquad
  \hat{S}^2 = \hat{S}_z^2
    + \frac{1}{2}\left(\hat{S}_+\hat{S}_- + \hat{S}_-\hat{S}_+\right),
  \qquad
  \hat{S}_{\pm} = \sum_p a^{\dagger}_{p\pm}a_{p\mp},
  \label{eq:4-pairingspin}
\end{equation}
so the Hamiltonian is block diagonal in $S$ and $S_z$.  We work in the
$S=0$ block.  There it is natural to introduce the pair operators
\begin{equation}
  \hat{P}^{+}_p = a^{\dagger}_{p+}a^{\dagger}_{p-},
  \qquad
  \hat{P}^{-}_p = a_{p-}a_{p+},
  \label{eq:4-pairoperators}
\end{equation}
in terms of which
\begin{equation}
  \hat{H} = \sum_{p\sigma}(p-1)a_{p\sigma}^{\dagger}a_{p\sigma}
    - \frac{1}{2}g\sum_{pq}\hat{P}^{+}_p\hat{P}^{-}_q .
  \label{eq:4-pairingpairform}
\end{equation}
The pair creation operators commute among themselves, $[\hat{P}^{+}_p,
\hat{P}^{+}_q]=0$, which is why they behave like bosons and why the model
describes a superfluid.  They are \emph{not} bosons, however:
$(\hat{P}^{+}_p)^2=0$, the Pauli principle asserting itself.

\paragraph{The seniority-zero space.}
Because the interaction only moves pairs, it never breaks one.  If we start
from a state in which every occupied level carries a full pair, we stay in
that space forever.  A basis state is then simply a choice of which $n$ of the
$L$ levels carry a pair, and the dimension is $\binom{L}{n}$ rather than
$\binom{2L}{2n}$.  For $L=4$ and $n=2$ this is $6$ instead of $28$; for $L=12$
and $n=6$ it is $924$ instead of $2704156$.  This is the space in which
Table~\ref{tab:lanczos} was computed.

In that basis the Hamiltonian is easy to write down: the diagonal is
$2\sum_{p\in\mathrm{occ}}(p-1) - \tfrac{1}{2}gn$, and any two configurations
differing by the position of a single pair are connected by $-\tfrac{1}{2}g$.

Table~\ref{tab:pairing} gives the exact ground-state energy for four particles
in four levels as a function of $g$, together with the energy of the
unperturbed configuration.  At $g=0$ the reference determinant is exact.  For
small $|g|$ the correlation energy grows quadratically, as second-order
perturbation theory requires, and for $g=1$ it has reached $36\%$ of the
reference energy.

\begin{table}[h]
\centering
\renewcommand{\arraystretch}{1.25}
\setlength{\tabcolsep}{12pt}
\begin{tabular}{rlll}
\toprule
$g$ & $E_{\mathrm{ref}}$ & $E_0$ (exact) & $E_{\mathrm{corr}}$\\
\midrule
$-1.00$ & $3.000000$ & $2.77987014$ & $-0.22012986$\\
$-0.50$ & $2.500000$ & $2.43688426$ & $-0.06311574$\\
$0.00$  & $2.000000$ & $2.00000000$ & $\phantom{-}0.00000000$\\
$0.50$  & $1.500000$ & $1.41677428$ & $-0.08322572$\\
$1.00$  & $1.000000$ & $0.63554847$ & $-0.36445153$\\
\bottomrule
\end{tabular}
\caption{The pairing model with four doubly degenerate levels and two pairs,
$\xi=1$.  $E_{\mathrm{ref}}$ is the energy of the configuration with both
pairs in the two lowest levels; $E_0$ is the exact ground-state energy in the
six-dimensional seniority-zero space.}
\label{tab:pairing}
\end{table}

\begin{notebox}
\textbf{Why this model matters.} The pairing model is the schematic version
of superfluidity and superconductivity, and it is the standard proving ground
for the hierarchy of many-body methods.  Because the exact answer is
available for any $g$, one can plot Hartree-Fock, second- and third-order
perturbation theory, coupled cluster with doubles, and full configuration
interaction on the same axes and watch each one fail at a different coupling.
The second midterm of FYS4480 does exactly this, and the exercises at the end
of this chapter set up the ingredients.
\end{notebox}

%----------------------------------------------------------------------
\section{Breaking the pairs: a particle-hole interaction}
\label{sec:pairingph}

\index{pairing model!with particle-hole interaction}
\index{seniority}
The pairing model is almost too well behaved.  Its interaction moves pairs and
does nothing else, so the seniority -- the number of particles that are not
part of a complete pair -- is a constant of the motion, and the whole
calculation collapses into the $\binom{L}{n}$-dimensional seniority-zero space.
That is a blessing when one wants exact answers cheaply, but it is a
liability when one wants to test methods, because an entire class of
excitations is simply absent.  In particular the \emph{particle-hole} channel,
in which a single particle is lifted across the Fermi level leaving one hole
behind, is empty: the pairing interaction has no matrix element between the
reference determinant and any one-particle-one-hole state.  The
Tamm-Dancoff and random-phase approximations of
Chapter~\ref{chap:mfapplications} are built entirely on that channel, and on
the pure pairing model they have nothing to say.

The cure is one extra term.  We keep the level scheme of
Eq.~(\ref{eq:4-pairingH}) and add an interaction that creates a pair while
removing two particles which need not be partners,
\begin{equation}
  \hat{H} = \hat{H}_0 + \hat{V}_{\mathrm{pair}} + \hat{V}_{\mathrm{ph}},
  \qquad
  \hat{V}_{\mathrm{ph}} = -\frac{f}{2}\sum_{pqr}
    \left(a^{\dagger}_{p+}a^{\dagger}_{p-}a_{q-}a_{r+}
      + \mathrm{h.c.}\right),
  \label{eq:4-Vph}
\end{equation}
with $\hat{H}_0$ and $\hat{V}_{\mathrm{pair}}$ exactly as in
Eq.~(\ref{eq:4-pairingH}).  The new strength $f$ is a second knob.  At $f=0$
we recover the pairing model and every number of Table~\ref{tab:pairing}; at
$g=0$ only pair breaking survives; and in between the two mechanisms compete.
This is the model of Chapter~\ref{chap:mfapplications}, introduced here so
that it sits alongside the others rather than arriving unannounced.

\paragraph{What the extra term actually does.}
The triple sum in Eq.~(\ref{eq:4-Vph}) hides three quite different processes,
and it is worth separating them.

\begin{itemize}
\item When $q=r$ the operator is
  $a^{\dagger}_{p+}a^{\dagger}_{p-}a_{q-}a_{q+}$, which is the pairing term
  again.  These pieces merely renormalise $g\to g+f$; they move whole pairs
  and change nothing structurally.
\item When $q=p$ and $r\ne p$ -- more transparently, in the Hermitian
  conjugate, $a^{\dagger}_{r+}a^{\dagger}_{p-}a_{p-}a_{p+}
  = a^{\dagger}_{r+}\hat{n}_{p-}a_{p+}$ -- a \emph{single} spin-up particle is
  carried from level $p$ to level $r$ while its former partner stays put.
  This is a one-particle-one-hole excitation, and it is the process the
  pairing interaction cannot perform.
\item When $q$, $r$ and $p$ are all different, two particles are removed from
  two different levels and a pair is deposited in a third.  The result is a
  two-particle-two-hole determinant with two unpaired particles.
\end{itemize}

Figure~\ref{fig:4-levels} shows the level scheme and one representative of
each excitation.  Panel~(a) is the reference determinant $\bket{c}$ with the
two lowest levels filled; panel~(b) is the only kind of excitation the pairing
term can produce, a whole pair moved across the Fermi level; panels~(c)
and~(d) are the one-particle-one-hole and pair-breaking
two-particle-two-hole determinants that only $\hat{V}_{\mathrm{ph}}$ can
reach.

\begin{figure}[htbp]
\centering
\newcommand{\levelpanel}[4]{%
  \begin{scope}[xshift=#2]
    \foreach \p in {1,2,3,4}{%
      \draw[thick,black!75] (-0.10,{0.74*(\p-1)}) -- (1.40,{0.74*(\p-1)});
      \coordinate (#1-\p-u) at (0.38,{0.74*(\p-1)});
      \coordinate (#1-\p-d) at (0.92,{0.74*(\p-1)});
      \fill[white] (#1-\p-u) circle (0.105);
      \fill[white] (#1-\p-d) circle (0.105);
      \draw (#1-\p-u) circle (0.105);
      \draw (#1-\p-d) circle (0.105);
    }
    \foreach \s in {#3}{\fill (#1-\s) circle (0.105);}
    \draw[dashed,black!55] (-0.28,{0.74*1.5}) -- (1.58,{0.74*1.5});
    \node[below,align=center,font=\footnotesize] at (0.65,-0.34) {#4};
  \end{scope}}
\begin{tikzpicture}[x=1cm,y=1cm,
   pairarrow/.style={-{Stealth[length=5pt,width=4pt]},thick,Blue3,
                     shorten >=2.5pt,shorten <=2.5pt},
   pharrow/.style={-{Stealth[length=5pt,width=4pt]},thick,Red3,
                   shorten >=2.5pt,shorten <=2.5pt}]

\levelpanel{A}{0cm}{1-u,1-d,2-u,2-d}%
  {(a) reference $\vert c\rangle$\\ seniority $0$}
\levelpanel{B}{2.55cm}{1-u,1-d,3-u,3-d}%
  {(b) $2p$--$2h$, $\hat{V}_{\mathrm{pair}}$\\ seniority $0$}
\levelpanel{C}{5.10cm}{1-u,1-d,2-d,3-u}%
  {(c) $1p$--$1h$, $\hat{V}_{\mathrm{ph}}$\\ seniority $2$}
\levelpanel{D}{7.65cm}{1-d,2-u,3-u,3-d}%
  {(d) $2p$--$2h$, $\hat{V}_{\mathrm{ph}}$\\ seniority $2$}

\draw[pairarrow] (B-2-u) to[bend left=50] (B-3-u);
\draw[pairarrow] (B-2-d) to[bend right=50] (B-3-d);
\draw[pharrow]   (C-2-u) to[bend left=50] (C-3-u);
\draw[pharrow]   (D-2-d) to[bend right=50] (D-3-d);
\draw[pharrow]   (D-1-u) to[out=155,in=205,looseness=1.1] (D-3-u);

\node[left,font=\footnotesize] at (-0.20,0)    {$p=1$};
\node[left,font=\footnotesize] at (-0.20,0.74) {$p=2$};
\node[left,font=\footnotesize] at (-0.20,1.48) {$p=3$};
\node[left,font=\footnotesize] at (-0.20,2.22) {$p=4$};
\node[right,font=\footnotesize] at (9.25,{0.74*1.5}) {$\varepsilon_F$};
\node[font=\footnotesize] at (0.38,2.62) {$+$};
\node[font=\footnotesize] at (0.92,2.62) {$-$};
\node[left,font=\footnotesize] at (-0.20,2.62) {$\sigma=$};
\end{tikzpicture}
\caption{Single-particle levels of the pairing model with $L=4$ doubly
degenerate levels, $\varepsilon_p=(p-1)\xi$, and $N=4$ particles.  Each level
carries two slots, $\sigma=+$ and $\sigma=-$; a filled circle is an occupied
spin-orbital and an open circle an empty one, so an open circle below
$\varepsilon_F$ is a hole.  (a) The reference determinant.  (b) A
two-particle-two-hole excitation produced by the pairing term
$\hat{V}_{\mathrm{pair}}$: a complete pair is carried from level $2$ to level
$3$ and the seniority stays zero.  (c) A one-particle-one-hole excitation and
(d) a pair-breaking two-particle-two-hole excitation, both of which require
the particle-hole term $\hat{V}_{\mathrm{ph}}$ of Eq.~(\ref{eq:4-Vph}) and
both of which have seniority two.  In the pure pairing model panels (c) and
(d) are unreachable from (a).}
\label{fig:4-levels}
\end{figure}

\paragraph{The price: the seniority-zero space is no longer closed.}
Because $\hat{V}_{\mathrm{ph}}$ breaks pairs, we can no longer work with the
six pair configurations of Section~\ref{sec:pairingmodel}.  What survives is
the weaker statement that both terms conserve $N_+$ and $N_-$ separately, so
$S_z=(N_+-N_-)/2$ is still a good quantum number.  For $L=4$ and $N=4$ the
full space has $\binom{8}{4}=70$ determinants and the balanced sector
$S_z=0$ has $36$ -- six times the seniority-zero dimension, and the price of
one extra term.

Table~\ref{tab:pairingphcoupling} makes the content of
Figure~\ref{fig:4-levels} quantitative.  It lists every determinant that the
Hamiltonian connects directly to the reference, classified by excitation rank
and seniority, together with the matrix element.  At $f=0$ there are four such
determinants and they are all of type~(b).  At $f\ne0$ two further families
appear, exactly as the figure suggests, and the seniority-zero matrix element
picks up the promised $g\to g+f$.

\begin{table}[h]
\centering
\renewcommand{\arraystretch}{1.25}
\setlength{\tabcolsep}{10pt}
\begin{tabular}{llrl}
\toprule
\textbf{Channel} & \textbf{Seniority} & \textbf{Number} &
$\bra{\Phi}\hat{H}\bket{c}$\\
\midrule
\multicolumn{4}{l}{\emph{$g=1$, $f=0$: the pairing model}}\\
$2p$--$2h$ & $0$ & $4$ & $-g/2$\\
\midrule
\multicolumn{4}{l}{\emph{$g=1$, $f=0.05$}}\\
$1p$--$1h$ & $2$ & $8$ & $\pm f/2$\\
$2p$--$2h$ & $0$ & $4$ & $-(g+f)/2$\\
$2p$--$2h$ & $2$ & $8$ & $\pm f/2$\\
\bottomrule
\end{tabular}
\caption{The determinants connected to the reference by one action of the
Hamiltonian, for $L=4$ and $N=4$ in the $S_z=0$ sector.  With $f=0$ only the
seniority-zero two-particle-two-hole states are reached, which is why the
particle-hole channel of the pairing model is empty.  Computed by
\texttt{models.py}.}
\label{tab:pairingphcoupling}
\end{table}

The effect on the ground state is shown in Table~\ref{tab:pairingph}.  The
first row is the pairing model, and reproduces the $g=1$ entry of
Table~\ref{tab:pairing} exactly -- a check worth making, since the two rows
are computed in spaces of dimension six and thirty-six respectively.  As $f$
grows the correlation energy grows much faster than the reference energy
falls, which is the signature of a second, competing mean field; the whole of
Chapter~\ref{chap:mfapplications} is about deciding which of the two to build
the fluctuations on.

\begin{table}[h]
\centering
\renewcommand{\arraystretch}{1.25}
\setlength{\tabcolsep}{12pt}
\begin{tabular}{rrlll}
\toprule
$g$ & $f$ & $E_{\mathrm{ref}}$ & $E_0$ (exact) & $E_{\mathrm{corr}}$\\
\midrule
$1.00$ & $0.00$ & $1.0000$ & $\phantom{-}0.63554847$ & $-0.36445153$\\
$1.00$ & $0.05$ & $0.9000$ & $\phantom{-}0.45058234$ & $-0.44941766$\\
$1.00$ & $0.20$ & $0.6000$ & $-0.18348455$ & $-0.78348455$\\
$1.00$ & $0.50$ & $0.0000$ & $-1.69173670$ & $-1.69173670$\\
\bottomrule
\end{tabular}
\caption{The pairing model with a particle-hole term, $L=4$, $N=4$, $\xi=1$,
diagonalised in the $36$-dimensional $S_z=0$ sector.  The first row is the
pairing model and agrees with Table~\ref{tab:pairing} to all digits shown.}
\label{tab:pairingph}
\end{table}

\begin{notebox}
\textbf{Why this model matters.} A model with one interaction tests one
approximation.  With two competing interactions and two strengths, the same
system can be driven from a regime where a pairing mean field is right to one
where a particle-hole mean field is, and one can watch each approximation fail
where it should.  That is precisely what Chapter~\ref{chap:mfapplications}
does with the Tamm-Dancoff approximation, the random-phase approximation, BCS
and the quasiparticle RPA, all measured against the exact diagonalisation of
this section.
\end{notebox}

%----------------------------------------------------------------------
\section{The Fermi-Hubbard model}
\label{sec:hubbard}

\index{Hubbard model}
The two models so far live in an abstract single-particle space with no
geometry.  The Fermi-Hubbard model introduces a lattice, and with it the
competition that dominates the physics of correlated electrons: kinetic energy
wants to delocalise the particles, the interaction wants to keep them apart.

Electrons hop between neighbouring sites of a lattice and pay an energy $U$
whenever two of them, necessarily of opposite spin, occupy the same site,
\begin{equation}
  \boxed{\;
  \hat{H} = -\sum_{\langle ij\rangle\sigma} t_{ij}
    \left(c^{\dagger}_{i\sigma}c_{j\sigma} + \mathrm{h.c.}\right)
  + U\sum_i \hat{n}_{i\uparrow}\hat{n}_{i\downarrow} \;}
  \label{eq:4-hubbardH}
\end{equation}
with $\hat{n}_{i\sigma}=c^{\dagger}_{i\sigma}c_{i\sigma}$.  This is again
Eq.~(\ref{eq:4-generalH}): the one-body part is the hopping matrix, and the
only non-vanishing two-body matrix element is the on-site one.

The single-particle basis is now the set of lattice sites, so the creation
operators carry a position label.  Fermionic antisymmetry is enforced by the
Jordan-Wigner string: $c^{\dagger}_i c_j$ acting on an occupation-number state
picks up the parity of the occupied sites strictly between $i$ and $j$, which
is exactly the phase computed by the bit-counting algorithm at the end of
Chapter~\ref{chap:secondquant}.

\paragraph{Symmetries and sector dimensions.}
The Hamiltonian conserves $N_{\uparrow}$ and $N_{\downarrow}$ separately, so
we may diagonalise in a fixed sector of dimension
\begin{equation}
  \dim = \binom{N}{N_{\uparrow}}\binom{N}{N_{\downarrow}} .
  \label{eq:4-hubbarddim}
\end{equation}
Better still, the hopping term acts on the two spin species independently, so
the sector factorises as a tensor product and the kinetic part of the
Hamiltonian is
\begin{equation}
  \bm{K}_{\uparrow}\otimes\bm{I}_{\downarrow}
  + \bm{I}_{\uparrow}\otimes\bm{K}_{\downarrow},
  \label{eq:4-hubbardkron}
\end{equation}
the tensor-product construction of Section~\ref{sec:tensor} used in earnest.
The interaction is diagonal in the occupation basis.  On a periodic chain
translational invariance gives a further reduction into momentum sectors.

\paragraph{From metal to Mott insulator.}
Table~\ref{tab:hubbard} follows the half-filled four-site ring as $U$ grows.
At $U=0$ the model is exact free fermions: the single-particle energies are
$\varepsilon_k=-2t\cos k$ with $k=0,\pm\pi/2,\pi$, and filling the two lowest
levels for each spin gives $E_0=-4t$, which the diagonalisation reproduces.
As $U$ increases, the double occupancy falls steadily -- the electrons avoid
each other -- and the ground-state energy rises towards zero.

\begin{table}[h]
\centering
\renewcommand{\arraystretch}{1.25}
\setlength{\tabcolsep}{12pt}
\begin{tabular}{rlll}
\toprule
$U/t$ & $E_0/t$ & $E_0/(Nt)$ & $\sum_i\langle n_{i\uparrow}n_{i\downarrow}\rangle$\\
\midrule
$0$  & $-4.00000000$ & $-1.000000$ & $1.000000$\\
$2$  & $-2.82842712$ & $-0.707107$ & $0.453082$\\
$4$  & $-2.10274848$ & $-0.525687$ & $0.287325$\\
$8$  & $-1.32023496$ & $-0.330059$ & $0.129992$\\
$16$ & $-0.72286432$ & $-0.180716$ & $0.042033$\\
\bottomrule
\end{tabular}
\caption{The half-filled four-site Hubbard ring, $N_{\uparrow}=N_{\downarrow}
=2$, sector dimension $36$.  The last column is the total double occupancy in
the ground state.}
\label{tab:hubbard}
\end{table}

\paragraph{The strong-coupling limit.}
When $U\gg t$ the doubly occupied states are pushed far up in energy and can
be eliminated perturbatively.  What remains is a model of localised spins,
\begin{equation}
  \hat{H}_{tJ} = \hat{P}_G\left[
    -\sum_{\langle ij\rangle\sigma}t_{ij}
      \left(c^{\dagger}_{i\sigma}c_{j\sigma}+\mathrm{h.c.}\right)
    \right]\hat{P}_G
  + \sum_{\langle ij\rangle}J_{ij}
    \left[\bm{S}_i\cdot\bm{S}_j
      - \tfrac14\hat{P}_G\hat{n}_i\hat{n}_j\hat{P}_G\right],
  \qquad J_{ij}=\frac{4t_{ij}^2}{U},
  \label{eq:4-tJ}
\end{equation}
where $\hat{P}_G$ is the Gutzwiller projector removing all doubly occupied
sites.  At half filling there is no empty site to hop into, the kinetic term
is frozen, and the $t$-$J$ model reduces to the Heisenberg antiferromagnet.

This is a prediction we can check.  For the four-site ring the Heisenberg part
contributes $-2J$ and the constant $-\tfrac14 n_i n_j$ term on four bonds
contributes $-J$, so $E_0\rightarrow-3J=-12t^2/U$.
Table~\ref{tab:hubbardstrong} shows the ratio $E_0/J$ converging to $-3$.

\begin{table}[h]
\centering
\renewcommand{\arraystretch}{1.25}
\setlength{\tabcolsep}{14pt}
\begin{tabular}{rll}
\toprule
$U/t$ & $E_0/t$ & $E_0/J$, \quad $J=4t^2/U$\\
\midrule
$8$   & $-1.3202349583$ & $-2.64047$\\
$16$  & $-0.7228643224$ & $-2.89146$\\
$32$  & $-0.3714104223$ & $-2.97128$\\
$64$  & $-0.1870445461$ & $-2.99271$\\
$128$ & $-0.0936928521$ & $-2.99817$\\
\bottomrule
\end{tabular}
\caption{The half-filled four-site Hubbard ring at strong coupling.  The ratio
converges to $-3$, confirming the mapping onto the Heisenberg model with
$J=4t^2/U$.}
\label{tab:hubbardstrong}
\end{table}

\begin{notebox}
\textbf{Why this model matters.} The Hubbard model is the minimal model of
strong correlation.  It is believed to contain the essential physics of the
cuprate superconductors, it is the standard benchmark for quantum Monte Carlo,
tensor-network and machine-learning wave functions, and it is among the first
targets for quantum simulation with cold atoms and with quantum computers.  It
is also the model in which the exponential wall of Section~\ref{sec:expdim}
bites hardest: the sector dimension~(\ref{eq:4-hubbarddim}) for a $4\times4$
lattice at half filling is already $1.66\times10^8$, and the two-dimensional
model remains unsolved.
\end{notebox}

%----------------------------------------------------------------------
\section{The Calogero-Sutherland model}
\label{sec:calogero}

\index{Calogero-Sutherland model}
The three models above are solvable because a symmetry cuts the Hilbert space
down to a manageable size.  The Calogero-Sutherland
model~\cite{calogero1971,sutherland1971} is solvable for a different and much
rarer reason: it is \emph{integrable}, and its ground state is known in closed
form for any number of particles.

The Calogero model describes $N$ particles on a line in a harmonic trap,
interacting through an inverse-square potential,
\begin{equation}
  \hat{H} = -\frac{\hbar^2}{2m}\sum_{i=1}^{N}
      \frac{\partial^2}{\partial x_i^2}
    + \frac{m\omega^2}{2}\sum_{i=1}^{N}x_i^2
    + \sum_{i<j}\frac{g}{(x_i-x_j)^2} .
  \label{eq:4-calogeroH}
\end{equation}
It is conventional to parametrise the coupling as
\begin{equation}
  g = \frac{\hbar^2}{m}\lambda(\lambda-1),
  \label{eq:4-calogerocoupling}
\end{equation}
so that $\lambda$ measures the interaction strength.  The Sutherland model is
the same system on a ring of circumference $L$, with the inverse square
replaced by the chord distance,
\begin{equation}
  \hat{H} = -\frac{\hbar^2}{2m}\sum_i\frac{\partial^2}{\partial x_i^2}
    + \frac{\hbar^2}{m}\left(\frac{\pi}{L}\right)^2\sum_{i<j}
      \frac{\lambda(\lambda-1)}{\sin^2\!\left(\pi(x_i-x_j)/L\right)} .
  \label{eq:4-sutherlandH}
\end{equation}

\paragraph{The exact ground state.}
In units $\hbar=m=1$ the ground state of Eq.~(\ref{eq:4-calogeroH}) is
\begin{equation}
  \boxed{\;
  \Psi_0(x_1,\dots,x_N)
   = \prod_{i<j}|x_i-x_j|^{\lambda}\,
     \exp\!\left(-\frac{\omega}{2}\sum_i x_i^2\right) \;}
  \label{eq:4-calogeroground}
\end{equation}
with energy
\begin{equation}
  E_0 = \omega\left[\frac{N}{2} + \frac{\lambda N(N-1)}{2}\right] .
  \label{eq:4-calogeroenergy}
\end{equation}
Exercise~\ref{ex:4-calogero} verifies this by direct substitution; the
inverse-square form of the potential is precisely what is needed for the
cross terms generated by the Laplacian to cancel.

The factor $\prod_{i<j}|x_i-x_j|^{\lambda}$ is a \emph{Jastrow factor}, and it
is the reason this model deserves a place in a book about many-body methods.
It is a correlated wave function of exactly the type that variational Monte
Carlo optimises numerically -- but here it is exact, for every $N$ and every
$\lambda$.  The model is therefore the sharpest available test of a Jastrow
ansatz: any method that cannot reproduce Eq.~(\ref{eq:4-calogeroground}) is
missing something structural.

Two limits are instructive.  At $\lambda=0$ the coupling $g$ vanishes and the
particles are free bosons in a trap, $E_0=\omega N/2$.  At $\lambda=1$ the
coupling vanishes again, but now
\begin{equation}
  \prod_{i<j}|x_i-x_j|
  \label{eq:4-vandermonde}
\end{equation}
is the absolute value of the Vandermonde determinant of
Section~\ref{sec:applications}, which is precisely the Slater determinant of
$N$ free fermions in a harmonic trap.  The Calogero model interpolates
continuously between bosons and fermions, and for non-integer $\lambda$ it
describes particles obeying \emph{fractional statistics}.

\paragraph{A numerical check.}
For $N=2$ the centre-of-mass and relative motions separate.  With
$x=x_1-x_2$ the relative Hamiltonian is
\begin{equation}
  \hat{h}_{\mathrm{rel}} = -\frac{d^2}{dx^2}
    + \frac{\omega^2}{4}x^2 + \frac{\lambda(\lambda-1)}{x^2},
  \label{eq:4-calogerorel}
\end{equation}
whose exact ground-state energy is $\omega(\lambda+\tfrac12)$.  Discretising
on a grid and diagonalising with the machinery of
Section~\ref{sec:schrodiag} reproduces this to better than $10^{-5}$ for
$\lambda$ between $1$ and $3$, as \texttt{models.py} confirms.  The
inverse-square barrier keeps the wave function away from the origin, so the
boundary condition $\Psi\rightarrow0$ as $x\rightarrow0$ is automatic for
$\lambda>0$.

\begin{notebox}
\textbf{Why this model matters.} The Calogero-Sutherland model sits at the
junction of several subjects.  Its ground state is a Laughlin-type wave
function, $\prod_{i<j}(z_i-z_j)^m e^{-\sum|z_i|^2/4}$, connecting it to the
fractional quantum Hall effect; its excitation spectrum is that of free
particles with fractional statistics; and the probability density
$|\Psi_0|^2\propto\prod_{i<j}|x_i-x_j|^{2\lambda}e^{-\omega\sum x_i^2}$ is
exactly the eigenvalue distribution of a random matrix ensemble, with
$\beta=2\lambda$ the Dyson index.  For our purposes it is above all a
benchmark with a known answer for arbitrary $N$ -- something none of the other
models in this chapter can offer.
\end{notebox}

%----------------------------------------------------------------------
\section{What the four models have in common}
\label{sec:modelssummary}

Table~\ref{tab:models} collects the four.  The pattern is worth stating
explicitly, because it is the pattern of the whole subject.

\begin{table}[h]
\centering
\renewcommand{\arraystretch}{1.3}
\setlength{\tabcolsep}{7pt}
\begin{tabular}{lllll}
\toprule
\textbf{Model} & \textbf{Single-particle space} & \textbf{Interaction} &
\textbf{Symmetry used} & \textbf{Exactly solvable}\\
\midrule
Lipkin   & two levels, degeneracy $\Omega$ & pair transfer, spin exchange
         & $SU(2)$ quasispin & yes, block by block\\
Pairing  & $L$ doubly degenerate levels & pair transfer
         & $S$, $S_z$, seniority & yes, small $L$\\
Hubbard  & lattice sites & on-site $U$
         & $N_{\uparrow}$, $N_{\downarrow}$, momentum & 1D only\\
Calogero & continuum, $N$ particles & inverse square
         & integrability & yes, any $N$\\
\bottomrule
\end{tabular}
\caption{The four models of this chapter.  Each is the general
Hamiltonian~(\ref{eq:4-generalH}) with a drastically restricted set of
two-body matrix elements, and in each case a symmetry reduces the dimension
of the problem.}
\label{tab:models}
\end{table}

In every case the route is the same.  Write the Hamiltonian in second
quantisation; identify the operators that commute with it; use them to block
diagonalise; and solve the blocks.  The first step is
Chapter~\ref{chap:secondquant}, the second is a commutator calculation of the
kind Wick's theorem makes routine, the third is the statement that a symmetry
gives a basis in which the matrix is block diagonal
(Section~\ref{sec:spectral}), and the fourth is the eigenvalue problem of
Sections~\ref{sec:householder} and~\ref{sec:lanczos}.

What separates these models from a realistic system is only the last step of
the reduction.  In a real nucleus or a real molecule the symmetries reduce the
dimension by orders of magnitude but not to five; the matrix is sparse but
enormous; and one must fall back on the approximate methods that occupy the
rest of this book.  The value of the models in this chapter is that they let
us watch those approximations succeed and fail against an answer we already
know.

%=============================================================================
\section{Exercises}
%=============================================================================

\subsection*{Exercise 1: the Lipkin model in the quasispin basis}
\label{ex:4-lipkin}
\addcontentsline{toc}{subsection}{Exercise 1: the Lipkin model}

\begin{enumerate}
  \item[a)] Show that the quasispin operators~(\ref{eq:4-quasispin}) obey the
  angular-momentum algebra~(\ref{eq:4-su2}), using only the anticommutation
  relations of Chapter~\ref{chap:secondquant}.
  \item[b)] Show that the Lipkin Hamiltonian~(\ref{eq:4-lipkinH}) can be
  rewritten as Eq.~(\ref{eq:4-lipkinquasispin}).
  \item[c)] Show that $[\hat{H},\hat{J}^2]=0$ and $[\hat{H},\hat{N}]=0$, and
  explain what this implies for the structure of the Hamiltonian matrix.
  \item[d)] Construct the five states $\bket{2,J_z}$ for four particles, and
  compute the matrix~(\ref{eq:4-lipkinmatrix1}).  Diagonalise it and compare
  with Eq.~(\ref{eq:4-lipkinground1}).
\end{enumerate}

\paragraph{Answer to a).}
Insert the definitions and use
$\{a_l,a_k\}=\{a^{\dagger}_l,a^{\dagger}_k\}=0$ and
$\{a^{\dagger}_l,a_k\}=\delta_{lk}$.  For the first commutator,
\begin{align}
  [\hat{J}_z,\hat{J}_{\pm}]
  &= \frac{1}{2}\sum_{pp'\sigma}\sigma\left(
     a_{p\sigma}^{\dagger}a_{p\sigma}a_{p'\pm}^{\dagger}a_{p'\mp}
   - a_{p'\pm}^{\dagger}a_{p'\mp}a_{p\sigma}^{\dagger}a_{p\sigma}\right).
  \label{eq:4-exa1}
\end{align}
Anticommute the operators of the second product into the order of the first.
Using $a_{p'\mp}a_{p\sigma}^{\dagger}=\delta_{pp'}\delta_{\mp\sigma}
- a_{p\sigma}^{\dagger}a_{p'\mp}$ and then
$a_{p\sigma}a_{p'\pm}^{\dagger}=\delta_{pp'}\delta_{\pm\sigma}
- a_{p'\pm}^{\dagger}a_{p\sigma}$, the four-operator terms cancel in pairs and
what survives is
\begin{equation}
  [\hat{J}_z,\hat{J}_{\pm}]
  = \frac{1}{2}\sum_{p\sigma}\sigma\left(
     \delta_{\pm\sigma}a_{p\sigma}^{\dagger}a_{p\mp}
   - \delta_{\mp\sigma}a_{p\pm}^{\dagger}a_{p\sigma}\right)
  = \pm\frac{1}{2}\sum_p\left(
     a_{p\pm}^{\dagger}a_{p\mp} + a_{p\pm}^{\dagger}a_{p\mp}\right)
  = \pm\hat{J}_{\pm},
  \label{eq:4-exa2}
\end{equation}
where the two Kronecker deltas each pick out one value of $\sigma$, and the
factor $\sigma=\pm1$ supplies the sign.  The same manipulation gives
$[\hat{J}_+,\hat{J}_-]=2\hat{J}_z$, and $[\hat{J}^2,\hat{J}_{\pm}]=
[\hat{J}^2,\hat{J}_z]=0$ then follows from the definition
$\hat{J}^2=\hat{J}_+\hat{J}_-+\hat{J}_z^2-\hat{J}_z$ exactly as for ordinary
angular momentum.

\paragraph{Answer to b).}
The one-body part is immediate: comparing
Eqs.~(\ref{eq:4-lipkinH0}) and~(\ref{eq:4-quasispin}) gives
$\hat{H}_0=\varepsilon\hat{J}_z$.

For $\hat{H}_1$, first anticommute the two annihilation operators,
$a_{p'-\sigma}a_{p-\sigma}=-a_{p-\sigma}a_{p'-\sigma}$, and then move
$a_{p-\sigma}$ past $a_{p'\sigma}^{\dagger}$ using
$a_{p-\sigma}a_{p'\sigma}^{\dagger}
=\delta_{pp'}\delta_{-\sigma\sigma}-a_{p'\sigma}^{\dagger}a_{p-\sigma}$.
Since $\delta_{-\sigma\sigma}=0$, the delta term drops out and
\begin{equation}
  \hat{H}_1 = \frac{1}{2}V\sum_{pp'\sigma}
    a_{p\sigma}^{\dagger}a_{p-\sigma}\,
    a_{p'\sigma}^{\dagger}a_{p'-\sigma} .
  \label{eq:4-exb1}
\end{equation}
The two factors are now independent sums over $p$ and $p'$.  Writing out
$\sigma=\pm$ separately,
\begin{equation}
  \hat{H}_1 = \frac{1}{2}V\left[
    \Big(\sum_p a_{p+}^{\dagger}a_{p-}\Big)
    \Big(\sum_{p'}a_{p'+}^{\dagger}a_{p'-}\Big)
  + \Big(\sum_p a_{p-}^{\dagger}a_{p+}\Big)
    \Big(\sum_{p'}a_{p'-}^{\dagger}a_{p'+}\Big)\right]
  = \frac{1}{2}V\left(\hat{J}_+^2+\hat{J}_-^2\right).
  \label{eq:4-exb2}
\end{equation}

For $\hat{H}_2$ the same two steps give
\begin{equation}
  \hat{H}_2 = \frac{1}{2}W\left(
    -\sum_{p\sigma}a_{p\sigma}^{\dagger}a_{p\sigma}
    + \sum_{pp'\sigma}a_{p\sigma}^{\dagger}a_{p-\sigma}
      a_{p'-\sigma}^{\dagger}a_{p'\sigma}\right),
  \label{eq:4-exb3}
\end{equation}
where this time $\delta_{-\sigma-\sigma}=1$ and the delta term survives as
$-\hat{N}$.  Splitting the remaining sum over $\sigma$ as before identifies it
as $\hat{J}_+\hat{J}_-+\hat{J}_-\hat{J}_+$, so
$\hat{H}_2=\tfrac12 W(-\hat{N}+\hat{J}_+\hat{J}_-+\hat{J}_-\hat{J}_+)$.
Adding the three pieces gives Eq.~(\ref{eq:4-lipkinquasispin}).

\paragraph{Answer to c).}
Written as in Eq.~(\ref{eq:4-lipkinquasispin}), $\hat{H}$ is built entirely
from $\hat{J}_z$, $\hat{J}_{\pm}$ and $\hat{N}$.  Since $\hat{J}^2$ commutes
with each of $\hat{J}_z$ and $\hat{J}_{\pm}$ by part a), and with $\hat{N}$
trivially, it commutes with $\hat{H}$.  For $\hat{N}$: the one-body term
conserves particle number, and every term of the interaction contains two
creation and two annihilation operators, so $[\hat{H},\hat{N}]=0$.

Both $J$ and $N$ are therefore good quantum numbers.  In a basis
$\bket{J,J_z}$ the Hamiltonian matrix is block diagonal, one block for each
$J$, of dimension $2J+1$.  Within a block, $\hat{J}_z$ is diagonal and
$\hat{J}_{\pm}^2$ connects $J_z$ to $J_z\pm2$, so the block splits once more
into even and odd $J_z$.  A problem of nominal dimension
$\binom{2\Omega}{N}$ has been reduced to a set of matrices of dimension at
most $N/2+1$.

\paragraph{Answer to d).}
Start from $\bket{2,-2}$, Eq.~(\ref{eq:4-lipkinlowest}), and apply
Eq.~(\ref{eq:4-ladder}).  With $J=2$ and $J_z=-2$ the factor is
$\sqrt{6-2}=2$, so
\begin{equation}
  \bket{2,-1} = \frac{1}{2}\hat{J}_+\bket{2,-2}
   = \frac{1}{2}\Big(
     a_{1+}^{\dagger}a_{2-}^{\dagger}a_{3-}^{\dagger}a_{4-}^{\dagger}
   + a_{1-}^{\dagger}a_{2+}^{\dagger}a_{3-}^{\dagger}a_{4-}^{\dagger}
   + a_{1-}^{\dagger}a_{2-}^{\dagger}a_{3+}^{\dagger}a_{4-}^{\dagger}
   + a_{1-}^{\dagger}a_{2-}^{\dagger}a_{3-}^{\dagger}a_{4+}^{\dagger}
     \Big)\bket{0},
  \label{eq:4-exd1}
\end{equation}
one term for each substate that can be flipped.  Applying $\hat{J}_+$ again
with the factor $\sqrt{6-0}=\sqrt6$ gives $\bket{2,0}$ as the normalised sum
of the six ways of flipping two of the four substates, and so on up to
$\bket{2,2}$ with all four spins up.

The matrix elements follow from Eq.~(\ref{eq:4-lipkinquasispin}).  The
diagonal is
\begin{equation}
  \bracket{2,J_z}{\hat{H}}{2,J_z}
   = \varepsilon J_z + \frac{W}{2}\Big(-N + 2\big[J(J+1)-J_z^2\big]\Big),
  \label{eq:4-exd2}
\end{equation}
since $\hat{J}_+\hat{J}_-+\hat{J}_-\hat{J}_+=2(\hat{J}^2-\hat{J}_z^2)$.  The
off-diagonal elements come from $\tfrac12 V(\hat{J}_+^2+\hat{J}_-^2)$,
\begin{equation}
  \bracket{2,J_z+2}{\hat{H}}{2,J_z}
   = \frac{V}{2}\sqrt{J(J+1)-J_z(J_z+1)}\,
     \sqrt{J(J+1)-(J_z+1)(J_z+2)} .
  \label{eq:4-exd3}
\end{equation}
With $\varepsilon=2$, $V=-1/3$, $W=-1/4$ and $N=4$ this gives
Eq.~(\ref{eq:4-lipkinmatrix1}); for instance the $(-2,0)$ element is
$\tfrac12 V\sqrt{4}\sqrt{6}=-\sqrt6/3$, and the $J_z=-1$ diagonal element is
$-2+\tfrac12(-\tfrac14)(-4+2[6-1])=-2-\tfrac34=-\tfrac{11}{4}$.
Diagonalising reproduces Eq.~(\ref{eq:4-lipkinground1}).

\subsection*{Exercise 2: the pairing model}
\label{ex:4-pairing}
\addcontentsline{toc}{subsection}{Exercise 2: the pairing model}

\begin{enumerate}
  \item[a)] Show that $\hat{H}_0$ and $\hat{V}$ of Eq.~(\ref{eq:4-pairingH})
  commute with $\hat{S}_z$ and $\hat{S}^2$ of Eq.~(\ref{eq:4-pairingspin}).
  \item[b)] Show that the Hamiltonian can be written in the pair
  form~(\ref{eq:4-pairingpairform}), and that the pair creation operators
  commute among themselves while $(\hat{P}^{+}_p)^2=0$.
  \item[c)] Construct the Hamiltonian matrix for four particles in the four
  lowest levels, with no broken pairs and $S=0$, and find its eigenvalues as a
  function of $g$.
\end{enumerate}

\paragraph{Answer to a).}
$\hat{H}_0$ is a sum of number operators $a^{\dagger}_{p\sigma}
a_{p\sigma}$, and $\hat{S}_z$ is a weighted sum of the same operators, so the
two commute term by term.  For $\hat{V}$, note that
$a^{\dagger}_{p+}a^{\dagger}_{p-}$ creates one particle with $\sigma=+1$ and
one with $\sigma=-1$, changing $S_z$ by $\tfrac12(+1)+\tfrac12(-1)=0$;
likewise $a_{q-}a_{q+}$ destroys such a pair.  Hence $[\hat{V},\hat{S}_z]=0$.

For $\hat{S}^2$ it is enough to show $[\hat{V},\hat{S}_{\pm}]=0$.  The
operator $\hat{S}_+=\sum_p a^{\dagger}_{p+}a_{p-}$ acting on a fully paired
level gives zero, because the state $\bket{p+}$ is already occupied and
$(a^{\dagger}_{p+})^2=0$; acting on an empty level it also gives zero.  Since
$\hat{V}$ maps the fully paired space into itself,
$\hat{S}_+\hat{V}$ and $\hat{V}\hat{S}_+$ both annihilate it, and the
commutator vanishes.  Both $S$ and $S_z$ are therefore good quantum numbers
and the Hamiltonian is block diagonal in them.

\paragraph{Answer to b).}
Insert the definitions~(\ref{eq:4-pairoperators}) into
Eq.~(\ref{eq:4-pairingH}):
\begin{equation}
  -\frac{1}{2}g\sum_{pq}a^{\dagger}_{p+}a^{\dagger}_{p-}a_{q-}a_{q+}
  = -\frac{1}{2}g\sum_{pq}\hat{P}^{+}_p\hat{P}^{-}_q ,
  \label{eq:4-exp1}
\end{equation}
which is Eq.~(\ref{eq:4-pairingpairform}).

For the commutator, $\hat{P}^{+}_p\hat{P}^{+}_q
= a^{\dagger}_{p+}a^{\dagger}_{p-}a^{\dagger}_{q+}a^{\dagger}_{q-}$.  Moving
$a^{\dagger}_{q+}a^{\dagger}_{q-}$ to the front requires four interchanges of
neighbouring creation operators, each contributing a minus sign, so the total
sign is $(-1)^4=+1$ and $[\hat{P}^{+}_p,\hat{P}^{+}_q]=0$.  The pair operators
therefore behave like bosons under exchange.

They are not bosons, however.  Since
$(a^{\dagger}_{p+})^2=0$ by the Pauli principle,
\begin{equation}
  \left(\hat{P}^{+}_p\right)^2
   = a^{\dagger}_{p+}a^{\dagger}_{p-}a^{\dagger}_{p+}a^{\dagger}_{p-}
   = -\left(a^{\dagger}_{p+}\right)^2
      \left(a^{\dagger}_{p-}\right)^2 = 0 .
  \label{eq:4-exp2}
\end{equation}
A level can hold one pair but not two.  This is the hard-core constraint that
distinguishes the pairing model from a gas of free bosons, and it is what
makes the model non-trivial.

\paragraph{Answer to c).}
With four levels and two pairs there are $\binom{4}{2}=6$ basis states, which
we label by the occupied levels: $(12)$, $(13)$, $(14)$, $(23)$, $(24)$,
$(34)$.  The diagonal element of a configuration is
$2\sum_{p\in\mathrm{occ}}(p-1) - g$, the last term coming from the $p=q$ terms
of the interaction with two pairs present.  Every pair of configurations
differing by the position of one pair is connected by $-g/2$; configurations
differing in both pairs are not connected.  Ordering the basis as above,
\begin{equation}
  \bm{H} =
  \begin{bmatrix}
    2-g & -g/2 & -g/2 & -g/2 & -g/2 & 0\\
    -g/2 & 4-g & -g/2 & -g/2 & 0 & -g/2\\
    -g/2 & -g/2 & 6-g & 0 & -g/2 & -g/2\\
    -g/2 & -g/2 & 0 & 6-g & -g/2 & -g/2\\
    -g/2 & 0 & -g/2 & -g/2 & 8-g & -g/2\\
    0 & -g/2 & -g/2 & -g/2 & -g/2 & 10-g
  \end{bmatrix}.
  \label{eq:4-exp3}
\end{equation}
Diagonalising for a range of $g$ gives Table~\ref{tab:pairing}.  Note that the
$(13)$ and $(23)$ configurations are degenerate at $g=0$ -- both have
unperturbed energy $6$ -- and that the correlation energy is negative for both
signs of $g$, since the interaction can always lower the energy by mixing
configurations.

\subsection*{Exercise 3: the Hubbard model}
\label{ex:4-hubbard}
\addcontentsline{toc}{subsection}{Exercise 3: the Hubbard model}

\begin{enumerate}
  \item[a)] Show that the Hubbard Hamiltonian~(\ref{eq:4-hubbardH}) conserves
  $N_{\uparrow}$ and $N_{\downarrow}$ separately, and hence $S_z$.
  \item[b)] Diagonalise the $U=0$ Hubbard ring analytically and verify the
  entry $E_0=-4t$ of Table~\ref{tab:hubbard}.
  \item[c)] Derive the effective spin Hamiltonian in the limit $U\gg t$ at
  half filling, and predict the ratio $E_0/J$ for the four-site ring.
\end{enumerate}

\paragraph{Answer to a).}
The hopping term contains $c^{\dagger}_{i\sigma}c_{j\sigma}$ with the
\emph{same} $\sigma$ on both operators, so it conserves the number of
particles of each spin separately.  The interaction is a product of number
operators and conserves everything.  Hence $[\hat{H},\hat{N}_{\uparrow}]
=[\hat{H},\hat{N}_{\downarrow}]=0$, and since
$\hat{S}_z=\tfrac12(\hat{N}_{\uparrow}-\hat{N}_{\downarrow})$, also
$[\hat{H},\hat{S}_z]=0$.  This is what allows the diagonalisation to be
carried out in a fixed sector of dimension~(\ref{eq:4-hubbarddim}).

\paragraph{Answer to b).}
At $U=0$ the Hamiltonian is a sum of independent one-body problems.  On a ring
of $N$ sites with periodic boundary conditions the hopping matrix is
diagonalised by plane waves $c^{\dagger}_{k\sigma}=N^{-1/2}\sum_j
e^{\imath k j}c^{\dagger}_{j\sigma}$ with $k=2\pi n/N$, giving
\begin{equation}
  \varepsilon_k = -2t\cos k .
  \label{eq:4-exh1}
\end{equation}
For $N=4$ the allowed momenta are $k=0,\pm\pi/2,\pi$ with energies
$-2t$, $0$, $0$, $+2t$.  Filling the two lowest for each spin gives
$-2t+0=-2t$ per spin species, hence $E_0=-4t$, which is the first line of
Table~\ref{tab:hubbard}.  The degeneracy at $k=\pm\pi/2$ means the $U=0$
ground state is not unique; switching on $U$ lifts the degeneracy.

\paragraph{Answer to c).}
At half filling and $U\gg t$ every site carries exactly one electron in the
ground state, and states with a doubly occupied site lie an energy $U$ higher.
Treat the hopping as a perturbation.  A single hop creates a doubly occupied
site and an empty one, costing $U$; hopping back restores a singly occupied
configuration.  Second-order perturbation theory therefore gives an energy
gain $-2t^2/U$ for each such virtual excursion, and the process is possible
only if the two spins involved are antiparallel, since the Pauli principle
forbids two electrons of the same spin on one site.

Collecting the terms gives the Heisenberg Hamiltonian
\begin{equation}
  \hat{H}_{\mathrm{eff}}
   = J\sum_{\langle ij\rangle}\left(\bm{S}_i\cdot\bm{S}_j
     - \frac{\hat{n}_i\hat{n}_j}{4}\right),
  \qquad J = \frac{4t^2}{U},
  \label{eq:4-exh2}
\end{equation}
which is the half-filled limit of the $t$-$J$ model, Eq.~(\ref{eq:4-tJ}).
For the four-site ring the Heisenberg ground-state energy is $-2J$, and the
constant $-\tfrac14 n_i n_j$ contributes $-\tfrac14$ per bond on four bonds,
that is $-J$.  Hence
\begin{equation}
  E_0 \longrightarrow -3J = -\frac{12t^2}{U},
  \qquad\text{so}\qquad \frac{E_0}{J}\longrightarrow-3,
  \label{eq:4-exh3}
\end{equation}
which is what Table~\ref{tab:hubbardstrong} confirms numerically.

\subsection*{Exercise 4: the Calogero ground state}
\label{ex:4-calogero}
\addcontentsline{toc}{subsection}{Exercise 4: the Calogero ground state}

\begin{enumerate}
  \item[a)] Verify by direct substitution that
  Eq.~(\ref{eq:4-calogeroground}) is an eigenstate of
  Eq.~(\ref{eq:4-calogeroH}) with energy~(\ref{eq:4-calogeroenergy}).
  \item[b)] Show that at $\lambda=1$ the ground state is the Slater
  determinant of $N$ free fermions in a harmonic trap.
  \item[c)] Separate the $N=2$ problem and verify the relative energy
  $\omega(\lambda+\tfrac12)$.
\end{enumerate}

\paragraph{Answer to a).}
Work in units $\hbar=m=1$ and write $\Psi_0=e^{-\Phi}$ with
\begin{equation}
  \Phi = \frac{\omega}{2}\sum_i x_i^2
       - \lambda\sum_{i<j}\ln|x_i-x_j| .
  \label{eq:4-exc1}
\end{equation}
For any such form,
\begin{equation}
  -\frac{1}{2}\sum_i\frac{\partial^2\Psi_0}{\partial x_i^2}
   = \frac{1}{2}\left[\sum_i\frac{\partial^2\Phi}{\partial x_i^2}
     - \sum_i\left(\frac{\partial\Phi}{\partial x_i}\right)^2\right]\Psi_0 .
  \label{eq:4-exc2}
\end{equation}
The derivatives are
\begin{equation}
  \frac{\partial\Phi}{\partial x_i} = \omega x_i
    - \lambda\sum_{j\ne i}\frac{1}{x_i-x_j},
  \qquad
  \frac{\partial^2\Phi}{\partial x_i^2} = \omega
    + \lambda\sum_{j\ne i}\frac{1}{(x_i-x_j)^2} .
  \label{eq:4-exc3}
\end{equation}
Squaring the first,
\begin{equation}
  \sum_i\left(\frac{\partial\Phi}{\partial x_i}\right)^2
   = \omega^2\sum_i x_i^2
   - 2\lambda\omega\sum_i\sum_{j\ne i}\frac{x_i}{x_i-x_j}
   + \lambda^2\sum_i\sum_{j\ne i}\sum_{k\ne i}
     \frac{1}{(x_i-x_j)(x_i-x_k)} .
  \label{eq:4-exc4}
\end{equation}
Two simplifications now do all the work.  In the second term, pairing $(i,j)$
with $(j,i)$ gives
$x_i/(x_i-x_j)+x_j/(x_j-x_i)=1$, so the double sum equals $N(N-1)/2$ times
two, that is $-2\lambda\omega\cdot N(N-1)/2$.  In the third term the pieces
with $j=k$ give $\lambda^2\sum_{i\ne j}(x_i-x_j)^{-2}$, while the pieces with
$j\ne k$ cancel identically by the three-term identity
\begin{equation}
  \frac{1}{(x_i-x_j)(x_i-x_k)}
  + \frac{1}{(x_j-x_i)(x_j-x_k)}
  + \frac{1}{(x_k-x_i)(x_k-x_j)} = 0 .
  \label{eq:4-exc5}
\end{equation}
Assembling everything, the inverse-square terms combine as
\begin{equation}
  \frac{1}{2}\left[\lambda\sum_{i\ne j}\frac{1}{(x_i-x_j)^2}
   - \lambda^2\sum_{i\ne j}\frac{1}{(x_i-x_j)^2}\right]
   = -\lambda(\lambda-1)\sum_{i<j}\frac{1}{(x_i-x_j)^2},
  \label{eq:4-exc6}
\end{equation}
which cancels the interaction term of the Hamiltonian exactly; the
$\omega^2 x_i^2$ terms cancel the trap; and what is left is the constant
\begin{equation}
  E_0 = \frac{1}{2}\left[N\omega + \lambda\omega N(N-1)\right]
      = \omega\left[\frac{N}{2}+\frac{\lambda N(N-1)}{2}\right],
  \label{eq:4-exc7}
\end{equation}
which is Eq.~(\ref{eq:4-calogeroenergy}).  The cancellation in
Eq.~(\ref{eq:4-exc6}) works only because the interaction falls off as the
inverse \emph{square}; this is the sense in which the model is fine-tuned to
be solvable.

\paragraph{Answer to b).}
At $\lambda=1$ the coupling $g=\lambda(\lambda-1)$ vanishes and the particles
are free, but the wave function
\begin{equation}
  \Psi_0 = \prod_{i<j}|x_i-x_j|\,e^{-\omega\sum_i x_i^2/2}
  \label{eq:4-exc8}
\end{equation}
is not the bosonic ground state.  Up to an overall sign it is the Vandermonde
determinant of Section~\ref{sec:applications},
\begin{equation}
  \prod_{i<j}(x_j-x_i) = \det\left[x_i^{\,j-1}\right],
  \label{eq:4-exc9}
\end{equation}
and multiplying each column by the Gaussian and taking suitable linear
combinations turns $x_i^{j-1}e^{-\omega x_i^2/2}$ into the Hermite functions
of Section~\ref{sec:schrodiag}.  The result is precisely the Slater
determinant of the $N$ lowest harmonic oscillator orbitals.  The energy check
is immediate: Eq.~(\ref{eq:4-calogeroenergy}) with $\lambda=1$ gives
$E_0=\omega N^2/2 = \omega\sum_{n=0}^{N-1}(n+\tfrac12)$, the sum of the $N$
lowest oscillator energies.

For non-integer $\lambda$ the wave function is neither symmetric nor
antisymmetric, and the particles obey fractional statistics interpolating
between the two.

\paragraph{Answer to c).}
Introduce $X=(x_1+x_2)/2$ and $x=x_1-x_2$.  The kinetic energy separates as
$-\tfrac14\partial_X^2-\partial_x^2$ and the trap as
$\omega^2 X^2 + \tfrac14\omega^2 x^2$, so with $\hbar=m=1$
\begin{equation}
  \hat{H} = \left[-\frac{1}{4}\frac{\partial^2}{\partial X^2}
     + \omega^2 X^2\right]
   + \left[-\frac{\partial^2}{\partial x^2}
     + \frac{\omega^2}{4}x^2
     + \frac{\lambda(\lambda-1)}{x^2}\right].
  \label{eq:4-exc10}
\end{equation}
The centre-of-mass part is a harmonic oscillator with ground-state energy
$\omega/2$.  The relative part is Eq.~(\ref{eq:4-calogerorel}); its ground
state is $x^{\lambda}e^{-\omega x^2/4}$ with energy $\omega(\lambda+\tfrac12)$,
as substitution confirms.  The total is
$\omega/2+\omega(\lambda+\tfrac12)=\omega(1+\lambda)$, which agrees with
Eq.~(\ref{eq:4-calogeroenergy}) for $N=2$.  Solving the relative problem
numerically on a grid, as in Section~\ref{sec:schrodiag}, reproduces
$\omega(\lambda+\tfrac12)$ to better than $10^{-5}$; see \texttt{models.py}.

\subsection*{Exercise 5: numerical explorations}
\addcontentsline{toc}{subsection}{Exercise 5: numerical explorations}

\begin{enumerate}
  \item[a)] Using \texttt{models.py}, plot the Lipkin ground-state energy and
  the weight of $\bket{2,-2}$ in the ground state as functions of $V$ at fixed
  $\varepsilon$ and $W$.  At roughly what value of $V/\varepsilon$ does the
  single-configuration picture break down?
  \item[b)] Repeat the pairing calculation of Table~\ref{tab:pairing} for six
  and eight levels.  How does the correlation energy at $g=1$ scale with the
  number of levels?
  \item[c)] Compute the double occupancy of the Hubbard ring as a function of
  $U/t$ and identify the crossover.  Compare a six-site ring with the
  four-site one.
  \item[d)] For the pairing model with twelve levels and six pairs, compare
  the exact ground-state energy from full diagonalisation with the Lanczos
  result of Table~\ref{tab:lanczos}, and with the Schmidt truncation of
  Table~\ref{tab:schmidt}.  How many Schmidt states are needed for the energy
  to be correct to six digits?
\end{enumerate}
